\aircrafttype{F+W Vampire}

The Eidgenössische Konstruktionswerkstätte (F+W) Vampire is a day fighter, fighter-bomber, and trainer. All are license-built versions of corresponding \typeref{de Havilland Vampire} versions.

\subsection{Versions}

\subsubsection{F+W Vampire FB.6}

The FB.6 is a license-built version of the de Havilland Vampire FB.6 day fighter and fighter-bomber. The “Late” version reflects the 1960 upgrade to install an ejection seat.

It served in the Swiss Flugwaffe from 1950 to 1990 (alongside de Havilland Vampire FB.6s), although from 1968 only as trainers.

\subsubsection{F+W Vampire T.55}

The T.55 is a license-built version of the de Havilland Vampire T.55 trainer, but with the ejection seats of the “Late” version.

It served in the Swiss Flugwaffe from 1955 to 1990.

\subsection{Armament and Stores}

All versions have a gun armament of four Hispano 20 mm cannons with 150 rounds per gun. 

A typical air-to-air load would be two 450L (100 imperial gallon) fuel tanks to increase endurance. 

A typical air-to-ground load would be eight RP-3 rockets and then either two 500 lb bombs or fuel tanks, depending on the mission radius. On short-range missions, two 1,000-lb bombs could be carried but without rockets.

\subsection{ADCs}

\begin{adclist}
    \adcitem{F+W Vampire FB.6}
    \adcitem{F+W Vampire FB.6 (Late)}
    \adcitem{F+W Vampire T.55}
\end{adclist}

\subsection{See Also}

\begin{typelist}
    \typeitem{de Havilland Vampire}
\end{typelist}